\chapter{数据要素：确权、定价与数据资产}

\begin{examguide}
\textbf{本章考情}：数据要素是 0258 各校专业课的最高频章节（\important{专硕·核心}），"数据二十条"三权分置、数据资产入表、数据要素市场建设是名词解释与论述题的必考素材；数据定价方法涉及估值计算题。

\textbf{核心考点}：（1）数据要素的经济学属性（非竞争性、部分排他性、衍生性）；（2）Samuelson 条件、公地悲剧与反公地悲剧的福利分析；（3）数据确权的"不可能三角"与三权分置的产权结构；（4）数据定价三方法（成本法、收益法、市场法）与产业链七环节；（5）数据资产入表的会计处理。

\textbf{本章主线}：数据的属性（7.1）→ 数据归谁（7.2）→ 数据怎么定价（7.3）→ 数据如何入表（7.4）。
\end{examguide}

\section{数据要素的经济学属性}

\subsection{从副产品到生产要素}

第 6 章的分析指出，平台的市场势力、算法的定价能力、监管的执法工具，最终都汇聚于同一个生产要素：数据。本章接着回答这一汇聚之后的三个制度问题：数据归谁（确权）、数据值多少（定价）、数据如何进入资产负债表（入表）。回答的起点是数据本身的经济学属性：数据长期被视为生产过程的副产品，交易记录、用户日志、设备读数，被存储却不被定价。数字技术的成熟使数据完成了要素化的跃迁：其一，存储与计算成本的指数下降（第 1 章摩尔定律）使海量数据的留存与分析在经济上可行；其二，机器学习使数据的价值转化（从原始记录到预测能力）成为可规模化的生产过程；其三，平台商业模式使数据的生成、聚合与应用形成完整闭环。\mn{中国数据制度时间线速记：2020 年数据列为第五大生产要素、2022 年"数据二十条"三权分置、2024 年入表新规施行，三节点分别对应要素地位、产权结构、会计入账（后两者见 7.2 节与 7.4 节）。}2020 年 3 月，中共中央、国务院印发《关于构建更加完善的要素市场化配置体制机制的意见》，正式将数据与土地、劳动力、资本、技术并列为第五大生产要素，数据的要素地位由此获得制度确认。

\begin{definition}[数据要素（data factor）]
数据要素是指以电子形式存在、通过计算与分析参与生产过程、能够独立或协同其他要素创造价值的数据资源。
\end{definition}

定义中值得拆解的关键词是"参与生产过程"。数据与信息的本质区别正在于此：信息是认知层面的概念（消除不确定性），数据要素是生产函数层面的概念（进入 $Y = AF(K, L, D)$），数据必须经过"资源化—资产化—资本化"的转化，实际参与价值创造，才成为要素；单纯存储而未被使用的数据只是"数据资源"，尚未成为"数据要素"。

\subsection{与传统要素的属性对照}

数据要素的经济学属性可以通过与传统要素的系统对照来把握。属性上的差异决定激励与制度的差异：传统要素的产权、定价与交易制度是围绕其竞争性与排他性设计的，数据一旦偏离这些属性，整套制度就必须重新设计。表~\ref{tab:datafactor} 给出七个维度的系统对照。

\begin{table}[htbp]
\centering
\setlength{\tabcolsep}{4pt}
\caption{数据要素与传统要素的属性对照}
\label{tab:datafactor}
\begin{tabular}{p{2.2cm}p{3.9cm}p{3.9cm}p{2.2cm}}
\hline
\textbf{属性} & \textbf{土地/劳动/资本} & \textbf{数据} & \textbf{理论后果} \\
\hline
竞争性 & 竞争性（一人使用排除他人使用） & \textbf{非竞争性}：可同时多主体使用 & 共享不减少单主体价值 \\
排他性 & 可排他（产权可执行） & \textbf{部分排他}：技术可加密，但易复制扩散 & 产权执行成本高 \\
边际成本 & 边际成本递增或不变（典型假设） & 复制边际成本趋近于零 & 供给激励特殊 \\
价值特征 & 价值确定性强 & 价值不确定性：依赖场景与算法 & 定价困难 \\
非均质性 & 同质，单位要素近似可替代 & \textbf{非均质性}：多源异构、质量参差、价值密度不一 & 数据质量评估与标准化困难 \\
衍生性 & 使用中逐渐消耗 & \textbf{衍生性}：聚合加工可衍生新数据、新价值 & 使用越多价值越大，与折旧逻辑相反 \\
边际收益 & 边际收益递减（给定其他要素） & 同质边际收益递减，多源异构带来新价值 & 数据优势来自多样性而非数量堆积 \\
\hline
\end{tabular}
\end{table}

四项属性的组合定义了数据要素的制度难题：非竞争性意味着\textbf{共享是有效率的}（同一数据多人使用不互相妨碍），部分排他性意味着\textbf{排他是可能的但不彻底的}（产权可以设定但执行困难），边际成本趋近于零意味着\textbf{传统的成本加成定价不适用}，价值不确定性意味着\textbf{定价必须依赖场景与预期}。数据要素市场的全部制度设计，确权、定价、交易，都是对这四项属性组合的回应。后两行的非均质性与衍生性同样不可忽视：非均质性使数据质量的评估与标准化成为交易的前置难题，衍生性则意味着数据的价值在使用中增长，与传统要素使用中消耗的折旧逻辑相反。数据要素的衍生性（derivability）指数据经聚合、加工、关联分析可以衍生出新的数据产品与信息价值，衍生的新数据本身又可再次衍生，价值链条没有天然的终点；这一属性使数据要素的长期价值难以在静态时点评估，也是数据区别于一般数字产品的关键。

其中"边际成本"一行的结论值得单独展开：复制边际成本趋近于零改变了企业的成本结构，而成本结构是定价与市场结构的起点。数据供给的成本中，采集、存储、治理的固定投入占比极高，复制与分发的边际成本可以忽略不计，这一结构对平均成本曲线的影响可以精确刻画。

\begin{derivation}{数据复制成本与平均成本曲线}{}
\textbf{复制成本趋近于零的平均成本曲线。}设数据产品的总成本由固定成本 $F$（采集、清洗、治理的投入）与复制成本构成，每复制一份的边际成本为 $c$，供给数量为 $Q$，则总成本与平均成本为
\[
TC = F + cQ, \qquad AC = \frac{F}{Q} + c.
\]
传统商品的边际成本随产量上升，平均成本呈 U 形：$AC$ 先因固定成本摊薄而下降，后因边际成本上升而回升。数据复制的边际成本趋近于零（$c \to 0$），平均成本退化为
\[
AC = \frac{F}{Q},
\]
一条随 $Q$ 增加持续下降的双曲线：固定成本被不断摊薄，而边际成本不再构成回升的力量，U 形结构消失。\important{复制边际成本趋近于零使数据供给呈现平均成本持续递减}；同时边际成本也趋近于零，按边际成本定价将使价格趋近于零，企业无法回收固定成本，而成本加成定价需要一个可锚定的平均成本最小值，这个锚在 $AC = F/Q$ 上同样不存在。
\end{derivation}

这一结论落在经验现象上有两个直接后果。其一，市场结构：数据与数字产品的供给存在自然的集中倾向，先发者以更大的销量摊薄固定成本，成本结构本身构成进入壁垒，与第 6 章网络效应与数据反馈回路的进入壁垒结论相互加强。其二，定价方式：数据产品的定价实践普遍放弃成本基础、转向价值基础，按使用效果、数据量与场景定价，这正是 7.3 节三种定价方法的背景。

\subsection{公共品视角：Samuelson 条件}

7.1.2 节的属性对照显示，数据兼具"非竞争"与"部分排他"两个特征，这在福利经济学中对应一类特殊商品：公共品。要回答"数据共享是过度还是不足"，首先需要公共品最优供给的基准条件。

\begin{definition}[公共品（public good）]
公共品指同时具有非竞争性（一人消费不减少他人可得数量）与非排他性（难以或不应排除他人消费）的商品，国防、路灯、基础研究皆属此类。数据具有非竞争性但\textbf{部分排他}（表~\ref{tab:datafactor}），属于介于私人品与纯公共品之间的俱乐部品/准公共品：共享是有效率的，但其供给激励与纯公共品不同。
\end{definition}

公共品理论的最优条件构成数据共享问题的理论基准。

\begin{derivation}{Samuelson 条件}{}
\textbf{Samuelson 条件（Samuelson condition）。}考虑一种公共品 $G$（数据共享服务）与一种私人品 $x$，社会中有两个消费者。公共品的生产成本 $C(G)$ 以私人品为单位计量：提供 $G$ 单位公共品需要放弃 $C(G)$ 单位私人品，故资源约束为
\[
x_1 + x_2 + C(G) \leq \bar x,
\]
其中 $\bar x$ 为社会的私人品禀赋总量。设社会偏好为功利主义形式 $W = u_1 + u_2$（社会总效用最大化；这一设定使加总结构最为清晰，更一般的福利函数会得到福利权重的推广形式），社会福利最大化问题为
\[
\max \; u_1(x_1, G) + u_2(x_2, G) \quad \text{s.t.} \quad x_1 + x_2 + C(G) \leq \bar x.
\]
构造拉格朗日函数
\[
\mathcal L = u_1(x_1, G) + u_2(x_2, G) + \lambda \left[ \bar x - x_1 - x_2 - C(G) \right],
\]
对私人品的一阶条件给出
\[
\frac{\partial u_1}{\partial x_1} = \lambda, \qquad \frac{\partial u_2}{\partial x_2} = \lambda,
\]
即两种商品的配置都使私人品边际效用等于影子价格 $\lambda$。对公共品的一阶条件为
\[
\frac{\partial u_1}{\partial G} + \frac{\partial u_2}{\partial G} = \lambda C'(G),
\]
左边是\textbf{加总}而非单人边际效用：公共品具有非竞争性，一单位 $G$ 同时进入两个消费者的效用函数，边际社会收益是两个消费者边际效用之和，这正是公共品与私人品的分水岭。将三个一阶条件合并，用 $\lambda = \partial u_1/\partial x_1 = \partial u_2/\partial x_2$ 消去影子价格，得
\[
\frac{\partial u_1/\partial G}{\partial u_1/\partial x_1} + \frac{\partial u_2/\partial G}{\partial u_2/\partial x_2} = C'(G),
\]
其中每一项 $\frac{\partial u_j/\partial G}{\partial u_j/\partial x_j}$ 是消费者 $j$ 的边际替代率（marginal rate of substitution，MRS）：保持效用不变时，消费者愿意为多获得一单位 $G$ 而放弃的私人品数量，即以私人品计的边际支付意愿；$C'(G)$ 是生产的边际转换率（marginal rate of transformation，MRT）：多生产一单位 $G$ 必须放弃的私人品数量，即边际成本。于是公共品的最优供给条件为
\[
{\color{blue} \sum MRS = MRT}.
\]
与私人品的最优条件 $MRS_j = MRT$（每个人的边际支付意愿各自等于边际成本）对照：一单位私人品只服务一个消费者，边际收益是单人的 $MRS$；一单位公共品同时服务所有消费者，边际收益是全体消费者 $MRS$ 的加总。公共品的"加总"性质使\textbf{个体的付费意愿之和}而非个体意愿本身成为定价基准。
\end{derivation}

\begin{figure}[htbp]
\centering
\begin{tikzpicture}[>=Stealth]
\draw[->, very thick] (0,0) -- (0,9.5);
\draw[->, very thick] (0,0) -- (10.5,0);
\draw[main, very thick] (0.5,2.85) -- (10,0) node[pos=0.6, above=3pt, font=\small, fill=white, inner sep=2pt] {$D_1$};
\draw[second, very thick] (0.5,5.5) -- (6,0) node[pos=0.3, above=3pt, font=\small, fill=white, inner sep=2pt] {$D_2$};
\draw[structurecolor, very thick] (0.5,8.35) -- (6,1.2) node[pos=0.5, above=3pt, font=\small, fill=white, inner sep=2pt] {社会需求 $S=\sum MRS$};
\draw[structurecolor, very thick] (6,1.2) -- (10,0);
\draw[very thick] (0.5,2) -- (9.8,2) node[right=3pt, font=\small] {$MC=MRT$};
\draw[dashed, very thick] (5.3846,0) -- (5.3846,9.5) node[above=3pt, font=\small, fill=white, inner sep=2pt] {$G^*$};
\draw[dashed, very thick] (3.333,0) -- (3.333,2);
\node[font=\small] at (3.333,-0.5) {私人 $G^p$};
\node[font=\small, rotate=90] at (-0.6,4.75) {边际支付意愿与边际成本};
\node[font=\small] at (5.25,-0.5) {$G$（数据共享量）};
\end{tikzpicture}
\caption{Samuelson 条件：社会需求的垂直加总}
\label{fig:samuelson}
\end{figure}

对数据而言，Samuelson 条件的经济含义可以归结为：按单个主体的支付意愿定价（私人品逻辑）将系统性低估数据的社会价值，因为数据共享的边际社会收益是全体潜在使用者的边际价值之和。图~\ref{fig:samuelson} 给出了这一条件的图形表达：社会需求曲线由两条个人需求曲线\textbf{垂直加总}而成，它与边际成本线的交点给出社会最优供给 $G^*$，而按个人支付意愿定价得到的私人供给 $G^p$ 落在 $G^*$ 的左侧，供给不足在图形上一目了然。\mn{Samuelson 条件大白话：公共品定价要把所有使用者的支付意愿加总再和成本比（$\sum MRS = MRT$），私人品只需单人意愿等于成本（$MRS_j = MRT$）。照搬私人品逻辑给数据定价，等于只按一个人的支付意愿算账，其他潜在使用者的价值被丢在账外。}私人品定价逻辑下的数据供给不足，正是"数据孤岛"（data silo）的福利论根源：每个主体只按自己的边际收益决定数据供给，忽视了对其他主体的溢出价值。这也为数据要素市场的公共政策提供了依据：数据共享设施（公共数据开放平台、行业数据空间）的公共投资具有效率基础，7.2 节公共数据的开放与授权运营正是这一理论的应用。

共享的福利分析还须区分两种对称的失灵。第一种是\textbf{公地悲剧}（tragedy of the commons）：共享设施无人承担治理投入，质量随使用退化，公共数据若全部免费开放而治理成本无人承担，就会停留在原始堆积状态，这是 7.2 节授权运营要解决的融资问题。第二种是\textbf{反公地悲剧}（tragedy of the anticommons）：数据的使用权被多个主体分别持有，任何整合使用都须取得全部持有者的同意，交易成本随权利碎片化急剧上升，资源反而无法使用。数据孤岛正是反公地悲剧的形态：每个平台都握有自己那份数据的排他权，跨主体整合分析在技术上可行，在产权上寸步难行。两种失灵方向相反，公地悲剧是过度开放无人治理，反公地悲剧是过度排他无法流通；数据要素市场的制度设计因此必须两头并进：以授权运营解决治理融资，以三权分置与流通促进化解权利碎片。

\subsection{数据的最优采集与使用}

数据的要素属性决定了其最优配置的独特结构。数据采集的成本结构与传统要素不同：实体要素的边际成本在超过一定规模后上升（资源稀缺），而数据的采集成本同样上升但原因不同：合规成本随数据量扩张、隐私阻力随采集深度上升、存储与治理成本随数据复杂度增长；数据使用的边际价值则随同质数据增加而递减（同一来源的数据边际信息增益下降），但随\textbf{多样性}增加而递增（异质数据揭示新维度）。

\begin{derivation}{最优数据量}{}
\textbf{最优数据量的决定。}设数据存量 $D$ 的使用价值为 $V(D)$，采集与治理成本为 $C(D)$。\important{私人最优}满足
\[
\max_D \; V(D) - C(D) \;\Rightarrow\; {\color{blue} V'(D^*) = C'(D^*)}.
\]
与传统要素决策的差异在于 $V'(D)$ 的形态：单一来源数据的边际价值随存量快速递减（$V'' < 0$，同质数据的边际信息增益下降），但多源异构数据的边际价值可以持续为正。二者的关系取决于 $V(D)$ 的自变量结构：$V'' < 0$ 描述的是\textbf{给定来源结构下}沿同一价值函数的边际变化，而新增一个异质数据来源相当于价值函数整体上移（$V(D)$ 变为更高的 $\tilde V(D)$），同质数据堆叠服从边际递减，多样性扩张带来函数本身的上移，两条性质并不矛盾，数据的价值函数取决于\textbf{来源结构}而非仅取决于总量。

福利层面的差异更为关键。单个企业的私人最优 $D^*$ 只考虑其自身的 $V'(D)$；当数据被多主体共享使用时，一单位数据同时进入所有使用者的价值函数（与公共品的非竞争性同构），社会净收益为 $\sum_j V_j(D) - C(D)$，其一阶条件为
\[
V'_{\text{社会}}(D) = \sum_j V_j'(D) > V'(D),
\]
即各使用者边际价值之和（与 7.1.3 节 Samuelson 条件的加总结构一致）。\important{私人最优的数据采集低于社会最优}。但这并不意味着应放任采集：数据使用同时产生隐私负外部性，即对数据主体的损害 $E(D)$（隐私泄露风险、大数据歧视、行为操纵的可能），完整的社会最优问题为 $\max_D \sum_j V_j(D) - C(D) - E(D)$，一阶条件为
\[
\sum_j V_j'(D) = C'(D) + E'(D).
\]
\important{数据政策的全部张力都在这一等式中}：共享价值、采集成本、隐私损害三者的边际权衡，社会最优的数据规模是使三者边际平衡的那一点。
\end{derivation}

\section{数据产权制度}

\subsection{产权理论的基础}

数据要素市场的第一个制度问题是确权：数据归谁、谁能使用、收益如何分配。第 6 章的分析表明，数据反馈回路会自我强化为竞争壁垒，而壁垒的制度基础正是数据权属的不清晰；反过来，权属不清也使数据的流通与交易难以开展。产权制度要回答的正是"初始权利如何配置、配置之后资源能否有效流动"，其经典分析框架由科斯定理设定，数据确权的讨论必须从它讲起。

\begin{definition}[科斯定理（Coase theorem）]
科斯定理断言：若交易成本为零，无论产权初始赋予何方，资源都可通过自愿交易达到有效配置，产权安排只影响收入分配而不影响效率；若交易成本为正，产权安排本身决定资源配置的效率。
\end{definition}

科斯定理的表述是抽象的，交易成本、有效配置都是概念，理解与答题都需要落到具体数值。下面用数据企业与数据主体的谈判算例给出科斯定理的数值证明：产权先赋予一方，再考察交易成本为零与为正两种情形下资源配置是否随之变化。

\begin{derivation}{科斯定理的数值证明}{}
\textbf{科斯定理在数据场景的数值证明。}设数据企业 $A$ 使用数据产生的收益为 $v = 100$ 万元，其数据使用对主体 $B$ 造成的隐私损害为 $h = 40$ 万元，净社会收益
\[
v - h = 100 - 40 = 60\ \text{万元} > 0,
\]
此时"允许使用数据"是社会最优的，产权配置的讨论才有意义；若 $v < h$，无论产权赋予谁数据都不应被使用，确权问题退化为单纯的禁令问题。\textbf{情形一：交易成本为零。}若产权赋予 $A$（有权使用数据），$B$ 无权阻止，$A$ 直接使用数据，净社会收益 $60$ 万元，结果有效。若产权赋予 $B$（有权阻止），$A$ 可向 $B$ 出价 $t$ 购买使用权，只要补偿额满足
\[
h < t < v, \qquad 40 < t < 100,
\]
$B$ 获得高于损害的补偿后同意交易，$A$ 使用数据，净社会收益仍为 $60$ 万元。\important{结论：交易成本为零时，无论产权赋予谁，数据都被有效配置}，产权只影响分配（$t$ 的大小决定 60 万元在 $A$、$B$ 之间的分割），不影响效率。\textbf{情形二：交易成本为正。}设谈判成本 $T = 80$ 万元（估值分歧带来的信息成本、合规审查成本、契约起草与执行成本），则
\[
T = 80 > v - h = 60,
\]
谈判成本超过了可供分割的净收益，补偿额无区间可谈：产权赋予 $B$ 时谈判失败，数据不被使用，社会损失 $60$ 万元；产权赋予 $A$ 时数据被直接使用，无需谈判，$B$ 承受损害，净社会收益 $60$ 万元。\important{结论：交易成本为正时，产权安排决定效率}，产权应赋予能够使交易成本最小化的一方，或在谈判不可能时将权利赋予"能使总价值最大化"的使用方式。
\end{derivation}

数据确权的特殊困难在于：同一数据上叠加着多重主体的利益（个人数据的主体、收集加工的企业、使用数据的第三方），传统"一物一权"的所有权框架难以覆盖。个人对其行为数据享有什么权利？企业对其加工形成的数据集享有什么权利？公共部门掌握的数据属于谁？\mn{科斯定理在数据场景的表述（论述题标准）：交易成本为零 → 确权只影响分配；交易成本为正 → 确权决定数据能否被利用。数据交易中交易成本天然高昂（价值不确定性、合规成本），确权制度的效率含义被放大。}三类问题的叠加使数据确权成为全球性难题：欧盟选择"个人数据保护"路径（GDPR 以控制权而非所有权为核心），美国倾向"市场自由+事后监管"，中国则走出一条结构性分置的道路。

这些困难的根源可以用\textbf{数据确权的"不可能三角"}统一刻画。

\begin{definition}[数据确权的"不可能三角"]
数据确权的"不可能三角"指数据确权制度的目标无法同时满足\textbf{充分流通使用}（发挥数据价值）、\textbf{生产激励}（保护数据生产者的投入回报）与\textbf{隐私保护}（控制数据的外部性损害）三者：数据的同时多主体使用特性使三者两两冲突，任何确权方案都只能优先满足其中之二。
\end{definition}

\begin{figure}[htbp]
\centering
\begin{tikzpicture}[>=Stealth,
    vtx/.style={rectangle, rounded corners=2mm, draw=main, ultra thick, fill=main!8, align=center, font=\small, minimum width=24mm, minimum height=10mm},
    lab/.style={font=\small, fill=white, inner sep=2pt}]
\draw[main, ultra thick] (0,0) -- (6,0) node[pos=0.5, above=4pt, lab] {排他与流通冲突};
\draw[main, ultra thick] (3,5.196) -- (6,0) node[pos=0.5, above=3pt, sloped, allow upside down=false, lab] {收集激励与隐私冲突};
\draw[main, ultra thick] (0,0) -- (3,5.196) node[pos=0.5, above=3pt, sloped, allow upside down=false, lab] {限制使用与流通冲突};
\node[vtx] at (0,0) {充分流通};
\node[vtx] at (6,0) {生产激励};
\node[vtx] at (3,5.196) {隐私保护};
\node[lab] at (3,1.73) {只能兼顾其二};
\end{tikzpicture}
\caption{数据确权的"不可能三角"}
\label{fig:trilemma}
\end{figure}

冲突的结构是：数据的非竞争性要求尽可能扩大使用（流通目标），但扩大使用摊薄了生产者的独占收益（激励目标受损）；保护生产者收益需要排他性授权，而排他与流通直接冲突；隐私保护要求限制数据的收集与使用范围，与流通目标同样相悖。图~\ref{fig:trilemma} 给出了三角的图形表达，三条边分别对应一组两两冲突，图形中心的结论是任何确权方案只能兼顾其二。\mn{不可能三角速记：流通、激励、隐私三者不能全要。考试先报三目标，再报两两冲突，最后落到"三权分置正是绕开三角的制度方案"。}中国方案正是以这个三角形为背景设计的：三权分置不再追求完整的排他所有权，而是把三个目标拆到不同权利上分别落实，流通经由数据产品经营权实现，激励经由加工使用权保障，隐私经由持有权与个人数据权益规则约束。确权思路从"一物一权"转向"权利束拆分"，是理解数据二十条制度逻辑的关键入口。

\subsection{数据二十条与三权分置}

2022 年 12 月 19 日，中共中央、国务院发布《关于构建数据基础制度更好发挥数据要素作用的意见》，因包含二十条政策举措而被称为"数据二十条"。文件确立了数据产权制度的中国方案：\important{数据资源持有权、数据加工使用权、数据产品经营权}三权分置。

\begin{definition}[数据产权三权分置（separation of three rights over data）]
数据产权三权分置是"数据二十条"确立的数据产权结构性分置方案：淡化数据所有权，将数据产权分解为\textbf{数据资源持有权}（对数据资源的控制与保管）、\textbf{数据加工使用权}（对数据进行加工、分析、利用）、\textbf{数据产品经营权}（对数据产品进行经营、流通并获取收益）三项权利，由不同主体分置行使。
\end{definition}

三权分置的制度逻辑是\textbf{以"使用权"为中心重构产权}。传统所有权框架下，数据确权的核心困难在于所有权的排他性与数据的非竞争性（共享使用不互相妨碍）相冲突；三权分置将产权束拆分：持有权解决"数据由谁控制保管"（对应数据安全与合规义务），加工使用权解决"谁可以加工分析"（对应数据价值创造），产品经营权解决"谁可以销售数据产品"（对应数据流通收益）。\mn{"数据二十条"三权分置记忆：持有权（谁能持有数据资源）、加工使用权（谁能加工使用）、产品经营权（谁能经营数据产品）。不设"数据所有权"概念，回避一物一权的确权困境。}三权的分离使"数据所有权"这一悬而未决的问题被搁置，而数据的流通使用不再依赖所有权归属的最终裁决，这是对科斯定理交易成本含义的制度回应：当确权成本极高时，绕过所有权、直接配置使用权，是降低交易成本的制度设计。

配套的分级分类体系构成三权分置的实施框架：\textbf{公共数据}（政府部门持有）以开放共享为原则、授权运营为补充；\textbf{企业数据}由企业依法依规持有、自主开发利用；\textbf{个人数据}以"知情同意"为基础，个人对其数据享有查阅、复制、更正、删除等权利（《个人信息保护法》）。三类数据的确权逻辑各不相同：公共数据强调公共品属性（开放即效率），企业数据强调投入回报（谁加工谁受益），个人数据强调人格权益优先（人格权高于财产权）。

公共数据的"授权运营"制度是三权分置在公共数据领域的具体实现。

\begin{definition}[公共数据授权运营（authorized operation of public data）]
公共数据授权运营指公共数据的持有者（政府部门）依法将公共数据的加工使用权授予符合条件的运营机构，由运营机构投入治理成本、形成数据产品并对外提供，政府通过收益分成或公共服务回购实现公共价值返还的制度安排。
\end{definition}

"数据二十条"对公共数据的基本立场是"推动用于公共治理、公益事业的公共数据有条件无偿使用，探索用于产业发展、行业发展的公共数据有条件有偿使用"。这一立场内含授权运营的经济逻辑。公共数据（政务数据、公共服务数据）的持有者是国家机关，其开放面临两难：直接免费开放实现公共品效率（7.1 节 Samuelson 条件的直接应用），但原始公共数据的价值密度低、需要加工治理才能使用，而加工治理的投入由谁承担？"授权运营"的机制设计正是对这一问题的回答：\important{以数据产品的经营权换取数据治理的投入}，政府将公共数据的加工使用权授予运营机构，运营机构投入治理成本、形成数据产品、在数据交易所流通获利，政府保留持有权并以收益分成或公共服务回购的方式实现公共价值。这一机制的经济学本质是把"公共品供给的融资问题"转化为"经营权的竞争性分配"：运营权的授予方式（特许、招标）决定治理投入的效率，收益分配规则决定公共价值的回收程度，授权运营的设计质量，决定了公共数据是走向"公地悲剧"（tragedy of the commons）还是"价值释放"。

\begin{misconception}
\textbf{误区："公共数据具有公共品属性，就应当全部免费开放。"}公共品的效率条件要求按边际成本定价，而非一律免费：数据共享的边际成本趋近于零，免费开放接近效率定价，但原始公共数据的价值密度低，须经清洗、脱敏、治理才能形成可用产品，治理投入本身构成供给成本。全部免费意味着治理投入无人承担，公共数据停留在"原始堆积"状态，这正是"公地悲剧"在数据领域的形态；授权运营以经营权换治理投入，正是为了避免这一结局。
\end{misconception}

\section{数据定价与交易}

\subsection{数据要素产业链}

数据从原始记录到最终价值变现，要经过一条完整的产业链。国家工业信息安全发展研究中心对数据要素产业链的划分是行业标准口径，以七个环节覆盖数据要素从产生到价值变现的全过程，这一口径也是单选题的直接命题点。

\begin{definition}[数据要素产业链（data industry chain）]
数据要素产业链指数据要素从产生到价值变现所经历的完整价值链条，包括\textbf{数据采集}（获取原始数据）、\textbf{数据存储}（保存与管理数据）、\textbf{数据加工}（清洗、标注、治理）、\textbf{数据流通}（共享、开放与交易）、\textbf{数据分析}（挖掘、建模与洞察）、\textbf{数据应用}（场景化使用创造价值）与\textbf{生态保障}（安全、评估、合规等配套服务）七个环节，其中生态保障贯穿其余六个环节。
\end{definition}

七个环节的价值分配并不均匀：采集与存储解决"数据从哪来"，加工与分析解决"数据变成什么"，应用解决"数据值多少"，流通则决定前六个环节创造的价值能否兑现为价格。\mn{产业链七环节速记：采、存、加、流、分、应、保，即采集、存储、加工、流通、分析、应用、生态保障；生态保障是"全程护航"环节，单选题常以"哪一项属于生态保障"的形式命题。}本章的定价方法（7.3.2 节）服务于加工与应用环节的价值评估，交易市场（7.3.3 节）是流通环节的制度形态，两者构成产业链上价值实现的关键一环。图~\ref{fig:chain} 给出了七环节的链条结构：前六个环节构成主链，生态保障以底座形态贯穿全程。

\begin{figure}[htbp]
\centering
\begin{tikzpicture}[>=Stealth,
    box/.style={rectangle, rounded corners=1.5mm, draw=main, thick, fill=main!8, align=center, font=\small, minimum width=16mm, minimum height=9mm},
    base/.style={rectangle, rounded corners=1.5mm, draw=structurecolor, thick, fill=structurecolor!8, align=center, font=\small, minimum width=62mm, minimum height=11mm},
    arr/.style={->, very thick, second}]
\node[box] (c1) at (0,0) {数据采集};
\node[box] (c2) at (2.6,0) {数据存储};
\node[box] (c3) at (5.2,0) {数据加工};
\node[box] (c4) at (0,-1.9) {数据流通};
\node[box] (c5) at (2.6,-1.9) {数据分析};
\node[box] (c6) at (5.2,-1.9) {数据应用};
\node[base] (eco) at (2.6,-3.9) {生态保障：安全、评估、合规等\\（配套服务，贯穿全程）};
\draw[arr] (c1) -- (c2);
\draw[arr] (c2) -- (c3);
\draw[arr, shorten <=2pt, shorten >=2pt] (5.2,-0.45) -- (0,-1.45);
\draw[arr] (c4) -- (c5);
\draw[arr] (c5) -- (c6);
\draw[->, very thick, structurecolor, shorten <=2pt, shorten >=2pt] (0.6,-3.35) -- (0.6,-2.35);
\draw[->, very thick, structurecolor, shorten <=2pt, shorten >=2pt] (2.6,-3.35) -- (2.6,-2.35);
\draw[->, very thick, structurecolor, shorten <=2pt, shorten >=2pt] (4.6,-3.35) -- (4.6,-2.35);
\end{tikzpicture}
\caption{数据要素产业链的七环节结构}
\label{fig:chain}
\end{figure}

\subsection{三种定价方法}

数据定价的困难源于其价值属性的特殊组合（表~\ref{tab:datafactor}）。资产评估的三种经典方法被引入数据定价，各有其适用边界。

\textbf{成本法}（cost approach）：以数据获取、存储、加工的历史成本为基础定价，数据资产在重新获取的情形下以重置成本（重新采集、治理同等数据的现实投入）为替代基础。适用于尚未形成明确收益预期、市场可比交易稀少的数据资源。其缺陷在数据场景被放大：数据的复制边际成本趋近于零，历史成本与数据的使用价值严重背离\mn{三方法适用速记：成本法看"花了多少"（低估价值）、收益法看"能赚多少"（依赖预期）、市场法看"别人卖多少"（依赖可比交易）。数据交易实践中三者常组合使用并相互校验。}：一份花费千万采集的数据可能毫无用途，一份顺手记录的行为日志可能价值连城。

\textbf{收益法}（income approach）：以数据预期产生的经济收益贴现定价，是理论最严密的方法。

\begin{derivation}{现金流折现模型}{}
\textbf{现金流折现模型。}设数据资产在 $t$ 期的预期净收益为 $CF_t$，折现率为 $r$，则数据价值为
\[
V = \sum_{t=1}^{\infty} \frac{CF_t}{(1+r)^t}.
\]
若净收益恒定（$CF_t = CF$），则
\[
V = \frac{CF}{r};
\]
若净收益以恒定增长率 $g$ 永续增长（$g < r$），则
\[
V = \frac{CF_1}{r - g},
\]
即戈登模型（Gordon growth model），它是现金流折现公式在恒定增长假设下的简化。折现率 $r$ 的构成：$r = r_f + \beta(r_m - r_f) + \text{数据特有风险溢价}$，其中 $r_f$ 为无风险利率，$r_m - r_f$ 为市场风险溢价，$\beta$ 为系统风险系数，前两项来自资本资产定价模型（capital asset pricing model，CAPM），第三项为数据资产的特有风险溢价（时效性风险、合规风险、技术风险）。
\end{derivation}

收益法的缺陷在于预期的高度主观性：\mn{折现率三层结构速记：$r = r_f + \beta(r_m - r_f) +$ 数据特有风险溢价；CAPM 认为资产的期望收益只应补偿不可分散的系统风险，$\beta$ 是数据资产收益与市场组合收益的联动程度。}数据的使用价值依赖场景与算法，同样的数据在不同使用者手中产生完全不同的现金流，估值结果因人而异。这既是收益法的局限，也是数据交易中"价值发现"机制（拍卖）存在的原因：当卖方无法确知各买方的场景价值时，让买方以竞价揭示各自的支付意愿，比卖方单方面估算更为可靠，这正是 7.3.3 节的主题。

\textbf{市场法}（market approach）：以可比数据交易的价格为基准调整定价，如按数据量、维度、时效、稀缺性等因素修正可比案例价格。其前提是存在活跃、透明、可比的交易市场，当前数据交易市场恰恰缺乏这一前提，市场法的运用受限于样本稀少。表~\ref{tab:pricing} 汇总了三种方法的定价基础、适用情形与主要缺陷。

\begin{table}[htbp]
\centering
\setlength{\tabcolsep}{4pt}
\caption{数据定价三种方法的比较}
\label{tab:pricing}
\begin{tabular}{p{1.3cm}p{4.5cm}p{3.1cm}p{3.1cm}}
\hline
\textbf{方法} & \textbf{定价基础} & \textbf{适用情形} & \textbf{主要缺陷} \\
\hline
成本法 & 获取、存储、加工的历史成本；重新获取情形下用重置成本 & 尚未形成明确收益预期、可比交易稀少的数据资源 & 成本与使用价值背离，低估高价值数据 \\
收益法 & 预期净收益贴现 $V=\sum_{t} CF_t/(1+r)^t$，永续增长下 $V=CF/(r-g)$ & 使用场景清晰、预期收益可合理估计的数据 & 预期高度主观，估值因人而异 \\
市场法 & 可比数据交易价格，按数据量、维度、时效修正 & 存在活跃、透明的可比交易市场 & 可比样本稀少，市场前提欠缺 \\
\hline
\end{tabular}
\end{table}

\subsection{数据拍卖与交易市场}

数据价值的高度场景依赖与事前不确定性，使拍卖成为理论上的最优定价机制：拍卖以竞争性出价实现价值发现，投标人各自按自己场景下的预期收益出价，价格由"最高场景价值"决定（第 2 章拍卖理论）。数据拍卖的具体形式中，维克瑞拍卖（次高价支付）因其"诚实出价"的占优策略性质而受到推崇，数据买方无需策略性地隐瞒自己的场景价值，机制设计简化了数据这种"价值主观性强"商品的交易。

\begin{derivation}{数据拍卖算例}{}
\textbf{数据拍卖的价值发现算例。}设一份消费行为数据集面向三个潜在买家拍卖，买家在其各自业务场景下的估价为 $v_1 = 80$、$v_2 = 50$、$v_3 = 30$（万元）。\textbf{维克瑞拍卖}（次高价支付）下，买家 1 按占优策略如实出价 $b_1 = 80$ 胜出，支付次高出价 $50$ 万元，数据所有者获得 $50$，买家 1 获得剩余 $80 - 50 = 30$。诚实出价为何是占优策略，可以一般性地证明。设买家 $i$ 的估价为 $v_i$，其余买家的最高出价为 $p$，$p$ 独立于买家 $i$ 的出价。若买家 $i$ 如实出价 $b_i = v_i$：当 $p < v_i$ 时胜出，支付 $p$，获得剩余 $v_i - p > 0$；当 $p > v_i$ 时落败，剩余为零，两种情形下都拿到了对自己有利的结果。若低报 $b_i < v_i$：当 $p$ 落在区间 $(b_i, v_i)$ 内时，如实出价本可胜出，低报却导致落败，放弃本可到手的剩余 $v_i - p > 0$；$p$ 在区间之外时结果与如实出价相同，低报只会变差。若高报 $b_i > v_i$：当 $p$ 落在区间 $(v_i, b_i)$ 内时，高报导致胜出，但支付 $p > v_i$，剩余 $v_i - p < 0$，而如实出价本可落败避免亏损；$p$ 在区间之外时结果相同，高报只会变差。\important{如实出价 $b_i = v_i$ 弱占优于一切其他出价}：偏离要么不改变结果，要么使自己的剩余下降。\textbf{一级价格密封拍卖}下，对称均衡出价函数为 $b(v) = \frac{n-1}{n}v$（第 2 章推导，均匀分布情形），买家 1 出价 $b(80) = \frac{2}{3}\times 80 \approx 53.3$ 万元胜出并支付，所有者获 $53.3$，买家 1 剩余 $26.7$。两种机制下，所有者收入的\textbf{期望}相同（收益等价定理，第 2 章），但单次实现值不同：本例中维克瑞机制成交价 $50$，低于一级价格机制的 $53.3$，期望等价不等于实现值等价。
\end{derivation}

这一算例揭示了数据交易的机制设计权衡：维克瑞拍卖对买家透明（诚实出价即可），但次高价支付的"价格发现"功能弱，成交价 $50$ 与最高估价 $80$ 之间留有可观的未提取剩余。对数据这种"同一数据在不同场景价值悬殊"的商品，机制选择的实质是\textbf{在提取剩余与降低投标策略复杂度之间的取舍}；实践中数据交易所常以"撮合+协商"起步、以拍卖作为稀缺高价值数据集的定价机制。

数据拍卖还有一个数据要素特有的机制问题：联合竞买。传统商品拍卖的输家在拍后无法使用该商品，买家之间只有竞争；数据则不同，赢家拍下后可以无限次复制与共享，输家完全可以通过与赢家联合的方式获得数据，买家之间同时存在联合的可能。这一性质把拍卖的竞争基础掏空了一块：若多个买家事先约定联合竞买、拍下后共享使用，竞价的激烈程度下降，成交价被压低，价格发现功能受损。卖方的应对要么是在机制上设防，例如设置保留价以对抗低价合谋；要么干脆放弃一次性的"拍断"式销售，转向按用户量计费的共享定价（订阅制、按次授权），把非竞争性数据从"卖一次"改为"卖多次"。联合竞买问题的根源仍是 7.1.2 节的非竞争性：数据可以同时服务多个主体，拍卖这类为竞争性商品设计的机制移植到数据上，必须随非竞争性做出改造。

中国数据交易市场的制度形态是\textbf{数据交易所}。

\begin{definition}[数据交易所（data exchange）]
数据交易所是依法设立、为数据交易提供集中场所与配套服务的市场基础设施，承担合规审查、登记确权、价格发现与生态培育等功能，与场外交易共同构成数据要素流通的市场体系。
\end{definition}

贵阳大数据交易所（2015 年成立，全国首家）、上海数据交易所（2021 年）、深圳数据交易所（2021 年）等数十家场内交易机构，加之外部的场外交易，构成"场内+场外"的双轨结构。\mn{数据交易所的功能定位（简答题）：合规撮合（入场数据须经合规审查）、登记确权（数据产品登记证书）、价格发现（撮合与竞价）、生态培育（数据商、第三方服务机构）。场内交易占比仍低的现实，源于数据交易的三大堵点。}数据流通的堵点可以概括为\important{三难}：\textbf{确权难}：数据来源合法性链条的证明成本高；\textbf{定价难}：价值场景依赖导致价格共识难以形成；\textbf{合规难}：个人信息保护与数据安全义务使交易双方承担高额合规成本。三难相互强化：确权不明则交易风险高，风险高则价格折让大，价格折让大则供给方不愿入场，"数据孤岛"由此形成。

\section{数据资产入表}

\subsection{制度背景}

财政部《企业数据资源相关会计处理暂行规定》自 \important{2024 年 1 月 1 日}起施行，\mn{数据资产入表制度锚点：财政部《企业数据资源相关会计处理暂行规定》，2024 年 1 月 1 日施行。确认路径：符合无形资产定义与确认条件的按无形资产核算，日常持有以备出售的按存货核算。}允许符合条件的数据资源确认为\textbf{无形资产}（intangible asset）或\textbf{存货}（inventory）并计入资产负债表。

\begin{definition}[数据资产（data asset）]
数据资产是指企业拥有或控制、能够带来预期经济利益、且成本能够可靠计量的数据资源；在《企业数据资源相关会计处理暂行规定》框架下，它具体体现为确认为无形资产或存货的数据资源。
\end{definition}

数据资产入表的经济意义在于打通了数据"资源化—资产化—资本化"链条的关键环节：数据一旦进入资产负债表，企业以数据资产进行质押融资、投资入股、转让处置便获得了会计基础，数据从\important{"表外资源"变为"表内资产"}，其资本化（数据资产融资）随之成为可能。

\begin{misconception}
\textbf{误区："企业拥有的数据都是资产。"}数据要成为会计意义上的资产，须同时满足"企业拥有或控制""很可能带来经济利益""成本能够可靠计量"三个条件。企业存储的原始日志、爬取的非结构化数据，若无法证明其经济利益很可能流入或无法可靠计量成本，只能留在表外；仅当数据资源经过治理形成可辨认、可控制、可计量的形态，才进入资产负债表。从"数据资源"到"数据资产"的距离，正是 7.1.1 节所说"资源化—资产化"的转化过程。
\end{misconception}

数据从原始记录到融资工具的转化，可以概括为一条三阶段链条（图~\ref{fig:valuechain}）：资源化把原始记录变为可用数据（采集、清洗、治理），资产化把可用数据变为表内资产（确权、定价、入表），资本化把表内资产变为融资工具（质押、证券化、投资）。三个阶段分别对应本章 7.1、7.2--7.4 与要素市场的进一步深化，链条上的每一步都以制度安排为前提。

\begin{figure}[htbp]
\centering
\begin{tikzpicture}[>=Stealth,
    stage/.style={rectangle, rounded corners=2.5mm, draw=main, ultra thick, fill=main!8, align=center, font=\small, minimum width=30mm, minimum height=14mm},
    note/.style={font=\small, color=structurecolor, align=center},
    arr/.style={->, very thick, second}]
\node[stage] (A) at (0,0) {数据资源化\\原始记录 $\to$ 可用数据};
\node[stage] (B) at (5.0,0) {数据资产化\\可用数据 $\to$ 表内资产};
\node[stage] (C) at (10.0,0) {数据资本化\\表内资产 $\to$ 融资工具};
\node[note] at (0,-1.35) {采集、清洗、治理\\（7.1 节）};
\node[note] at (5.0,-1.35) {确权、定价、入表\\（7.2--7.4 节）};
\node[note] at (10.0,-1.35) {质押、证券化、投资\\（要素市场深化）};
\begin{scope}[on background layer]
\draw[arr, shorten <=-2pt, shorten >=-2pt] (A) -- (B);
\draw[arr, shorten <=-2pt, shorten >=-2pt] (B) -- (C);
\end{scope}
\end{tikzpicture}
\caption{数据价值化的三阶段链条}
\label{fig:valuechain}
\end{figure}

\subsection{确认条件与计量}

《暂行规定》并未新设一类资产，而是沿用《企业会计准则》的既有框架：数据资源能否入表，取决于它是否满足无形资产或存货的定义与确认条件，入表是在既有资产类别中"嵌套"数据资源。

数据资源确认为无形资产须同时满足：\textbf{可辨认性}（能够单独或与相关合同、资产一起识别）、\textbf{企业对其拥有控制权}（合法持有并有权获得其产生的经济利益）、\textbf{经济利益很可能流入}、\textbf{成本能够可靠计量}。这四项条件与《企业会计准则第 6 号——无形资产》的确认条件一致。\mn{入表确认四条件速记：可辨认、有控制、经济利益很可能流入、成本能够可靠计量，四者缺一不可；按用途归账：自用进无形资产，待售进存货。}确认为存货的条件类似，但以"持有以备出售"为目的（数据产品服务商的数据集库存）。

\begin{derivation}{摊销算例}{}
\textbf{摊销的数值算例。}设企业以 1000 万元外购一项数据资源，确认为无形资产，预计使用寿命 5 年，无残值，采用直线法摊销，则年摊销额为
\[
\text{年摊销额} = \frac{1000}{5} = 200\ \text{万元}.
\]
每期摊销的分录为：借记管理费用（或其他受益科目）200 万元，贷记累计摊销 200 万元；无形资产账面价值每年递减 200 万元。若该数据资源前期的经济价值显著高于后期（数据新鲜度随技术迭代衰减），采用年数总和法加速摊销更能实现费用与收益的配比。使用年限总和为
\[
5 + 4 + 3 + 2 + 1 = 15,
\]
第一年摊销额 $1000 \times \frac{5}{15} \approx 333.3$ 万元，第二年摊销额 $1000 \times \frac{4}{15} \approx 266.7$ 万元，此后逐年递减。对比直线法的固定 $200$ 万元，加速摊销把费用前移：前期利润更低、后期利润更高，两条路径的总摊销额同为 $1000$ 万元，差异仅在费用的时间分布。
\end{derivation}

数据资产摊销的独特难点在于\textbf{时效性}：数据价值的衰减通常快于传统无形资产。第 6 章的数据反馈回路模型已刻画了这一特征：数据转化的边际贡献递减（$f'' < 0$，推荐精度的边际提升随存量下降），任何数据优势都存在天然的饱和点。会计上对这一衰减的回应，正是估计使用寿命时充分考虑技术迭代与数据新鲜度，必要时采用加速摊销；使用寿命的估计由此成为数据资产摊销中最富争议的技术环节。

存货路径的计量同样有自己的数值逻辑。数据存货通常以外购方式获得，入账成本按历史成本可靠计量；但数据存货的价值衰减快于实物存货，期末计量适用与一般存货相同的谨慎规则：成本与可变现净值孰低。

\begin{derivation}{存货计量的数值算例}{}
\textbf{存货成本计量的数值算例。}设企业为出售而购入一项数据产品，采购成本 $100$ 万元，期末该数据产品因更新版本发布、同类数据集降价而贬值，预计售价降至 $60$ 万元，估计销售费用 $5$ 万元，则可变现净值（net realizable value）为
\[
\text{可变现净值} = 60 - 5 = 55\ \text{万元}.
\]
可变现净值 $55$ 万元低于成本 $100$ 万元，期末计提存货跌价准备
\[
100 - 55 = 45\ \text{万元},
\]
资产负债表上该存货按 $55$ 万元列报。若后续市价回升，可变现净值高于成本，存货仍按成本计量，减值可在原计提范围内转回，列报金额不超过成本。\important{数据存货按成本与可变现净值孰低计量}：贬值时及时减记，价值回升时不确认增值，谨慎性原则在数据存货上的表达与一般存货一致。
\end{derivation}

对比两条计量路径：无形资产路径的费用化方式是摊销（费用的时间分布问题），存货路径的减值方式是跌价准备（账面价值的时点校准问题），两者的共同底色都是数据价值的快速衰减。7.4.1 节的图~\ref{fig:valuechain} 显示资产化是链条的中段，而计量的质量直接决定资本化能否顺利衔接：账面价值可信，质押融资才可行。

\subsection{入表的影响}

数据资产入表的影响可以分三个层面理解。\textbf{企业层面}：资产负债表扩张（资产增加），利润表压力分化（外购数据资本化使当期费用递延、利润改善；内部开发的数据适用与内部研发无形资产相同的规则：研究阶段，即探索性工作、成果能否形成资产尚不确定的阶段，支出费用化；开发阶段，即有明确应用目标、具备资本化条件的阶段，支出方可资本化。研究阶段与开发阶段的划分是内部数据资产入表的实务考点）。杠杆效应随之显现：资产端的扩张改善资产负债率，更重要的是，入表后的数据资产具备质押融资的会计基础，企业以数据资产质押取得贷款，资产与负债联动扩张，融资能力由此提升。\textbf{市场层面}：数据资产的会计确认降低了数据交易的信息不对称，交易对手可以在报表中核查资产真实性，为数据质押融资提供估值基准。信息披露方面，《暂行规定》要求企业在财务报表附注中披露数据资源的相关信息，包括确认的类别、账面原值、累计摊销或跌价准备、摊销方法与使用寿命等，数据资产的构成与估值假设由此进入公众视野，披露的充分性是数据资产会计信息质量的关键。\textbf{宏观层面}：数据资本化的统计基础使数据要素对经济增长的贡献可核算化，为第 8 章"数据要素增长核算"提供了微观会计支撑。争议同样存在：数据价值的场景依赖性与表内资产的确定性要求之间存在张力，过度入表可能虚增资产，审慎性要求"证据充分、计量可靠"的数据方可入表。

\begin{chapterbox}
\textbf{本章考点清单}

\begin{enumerate}
\item 数据要素的属性（非竞争性、部分排他性、衍生性等七个维度）与制度难题 \important{【专硕·核心】}
\item Samuelson 条件 $\sum MRS = MRT$ 的推导、公地悲剧与反公地悲剧的福利分析 \important{【专硕·核心】}
\item 科斯定理的数据场景数值证明（交易成本为零/为正两情形）\important{【专硕·核心】}
\item 数据确权的"不可能三角"（流通、激励、隐私）与三权分置的制度回应 \important{【专硕·核心】}
\item "数据二十条"三权分置：持有权、加工使用权、产品经营权 \important{【专硕·核心】}
\item 公共数据/企业数据/个人数据的分类确权逻辑与公共数据授权运营 \important{【专硕·核心】}
\item 数据要素产业链七环节（采集、存储、加工、流通、分析、应用、生态保障）\important{【专硕·核心】}
\item 成本法、收益法（$V = CF/(r-g)$）、市场法的适用边界 \important{【专硕·核心】}
\item 数据拍卖的价值发现、维克瑞拍卖诚实出价证明与联合竞买问题 \important{【专硕·扩展】}
\item 数据交易市场结构与三大堵点（确权难、定价难、合规难）\important{【专硕·核心】}
\item 数据资产入表：确认条件、摊销与存货计量算例、三层面影响与信息披露 \important{【专硕·核心】}
\end{enumerate}
\end{chapterbox}

本章的分析始终在微观制度层面展开：确权降低交易成本，定价实现价值发现，入表提供会计基础。但数据要素的全部制度投入，最终要接受一个宏观尺度的检验：它是否真的改变了增长的路径？这一检验的标尺是增长核算，而数字技术在那里留下了一个著名的谜题。
